% TOP_PROBLEM: {"id": 3336, "title": "Termination of the sixth Goodstein Sequence", "category_id": 18, "category": {"id": 18, "name": "logic", "display_name": "Logic", "description": "Problems in mathematical logic, model theory, proof theory, and finite model theory.", "slug": "logic", "order_index": 18, "created_at": "2026-07-31T15:26:25.671Z"}, "status": "open"}
% FIRST_POSTED: null
% ENTRY_KIND: statement_only
% Imported from: batch03.tex; UnsolvedMath ID 3336
\documentclass[11pt]{article}
\usepackage[margin=25mm]{geometry}
\usepackage{fontspec}
\setmainfont{Latin Modern Roman}
\usepackage{amsmath,amssymb,mathtools}
\usepackage{unicode-math}
\setmathfont{Latin Modern Math}
\usepackage{xurl,hyperref}
\hypersetup{colorlinks=true,urlcolor=blue}
\setcounter{secnumdepth}{0}
\setlength{\parindent}{0pt}
\setlength{\parskip}{5pt}
\setlength{\emergencystretch}{3em}
\newcommand{\sref}[2]{\hyperlink{source:#1:#2}{[S#2]}}
\newcommand{\cataloguescope}[1]{\paragraph{Recorded target.}#1}
\begin{document}
% UnsolvedMath ID 3336; source catalogue code OPG-2379
\setcounter{section}{3335}
\section{Termination of the sixth Goodstein Sequence}\label{3336}
{\small\sffamily UnsolvedMath ID 3336 · Logic}
\subsection{Definitions and mathematical statement}

\paragraph{Shared notation.}$\mathbb N=\{1,2,\ldots\}$, and $\mathbb Z,\mathbb Q,\mathbb R,\mathbb C$ denote the integers, rationals, reals and complex numbers. Any other conventions are specified locally.


For $b\geq2$, the hereditary base-$b$ representation of a nonnegative integer writes its ordinary base-$b$ expansion $\sum c_i b^{e_i}$, and recursively writes each exponent $e_i$ in hereditary base $b$ as well. Let $H_{b\to b+1}(a)$ be the integer obtained by replacing every occurrence of the base $b$ in this representation by $b+1$, leaving digits unchanged. Define $g_0=6$ and, while $g_t>0$,\[ g_{t+1}=H_{t+2\to t+3}(g_t)-1. \] The termination step count is $T=\min\{t\geq0:g_t=0\}$. This counts transitions, so under the source's one-based term numbering the first zero is term $T+1$.
\paragraph{Statement recorded in the supplied source.}
Determine the exact termination step count $T$ for the sixth Goodstein sequence, starting from $6$ and hereditary base $2$. \sref{3336}{2}
\subsection{Short English statement}
How many successive hereditary base changes followed by subtracting one are needed for the Goodstein sequence starting at six to first reach zero?
\subsection{Sources}
\begin{description}
\item[\hypertarget{source:3336:1}{S1}] ULAM.AI and the UnsolvedMath Contributors, UnsolvedMath: A Curated Collection of Open Mathematics Problems, record 3336 (OPG-2379). CC BY 4.0. \url{https://huggingface.co/datasets/ulamai/UnsolvedMath}.
\item[\hypertarget{source:3336:2}{S2}] Open Problem Garden, Termination of the sixth Goodstein Sequence, node 2379, original problem page. Original question and full discussion. \url{https://www.openproblemgarden.org/op/termination_of_the_sixth_goodstein_sequence}.

\end{description}
\label{end:3336}
\end{document}
